% =============================================================================
% Apêndices C e D — Questionários da avaliação com profissionais
% Incluído por estrutura/pos-textuais/apendices.tex, dentro de apendicesenv.
% Versão em texto simples, para montar no Google Forms: docs/tcc/formularios.md
% =============================================================================

\providecommand{\caixa}{$\square$}
\providecommand{\opcao}[1]{\caixa~#1\quad}

\chapter{Questionário de adesão}
\label{ape:questionario-adesao}

Instrumento aplicado em formulário eletrônico a bibliotecários e estudantes de Biblioteconomia, sem necessidade de uso prévio da ferramenta. Tempo estimado de resposta: 8 a 10 minutos.

\section*{Parte 1 --- Termo de consentimento livre e esclarecido}

Você está sendo convidado(a) a participar da pesquisa ``Cataloga+: sistema de catalogação assistida de livros'', desenvolvida como Trabalho de Conclusão de Curso do curso de Tecnologia em Sistemas para Internet do Instituto Federal de Brasília. O objetivo deste questionário é verificar se as funcionalidades de uma ferramenta de apoio à catalogação descritiva atendem às necessidades da prática bibliotecária. A participação é voluntária e anônima: nenhuma informação que permita identificá-lo(a) será solicitada, e você pode desistir a qualquer momento, sem nenhum prejuízo. Os dados serão utilizados exclusivamente para fins acadêmicos e apresentados apenas de forma agregada. Em caso de dúvidas, entre em contato com o pesquisador pelo e-mail [inserir e-mail institucional].

\opcao{Li e concordo em participar da pesquisa.}

\section*{Parte 2 --- Perfil do respondente}

\begin{enumerate}
    \item Qual é a sua formação?\\
    \opcao{Graduação em Biblioteconomia em andamento}\opcao{Graduação em Biblioteconomia concluída}\opcao{Especialização}\opcao{Mestrado ou doutorado}\opcao{Outra}
    \item Há quanto tempo você atua com catalogação?\\
    \opcao{Não atuo}\opcao{Menos de 1 ano}\opcao{1 a 5 anos}\opcao{6 a 10 anos}\opcao{Mais de 10 anos}
    \item Em que tipo de biblioteca você atua?\\
    \opcao{Pública}\opcao{Escolar}\opcao{Universitária}\opcao{Especializada ou corporativa}\opcao{Comunitária}\opcao{Não atuo}\opcao{Outra}
    \item Qual é o tamanho aproximado do acervo?\\
    \opcao{Até 5 mil exemplares}\opcao{5 a 20 mil}\opcao{20 a 100 mil}\opcao{Mais de 100 mil}\opcao{Não sei}
    \item Qual sistema a sua biblioteca utiliza para catalogação? (pode marcar mais de uma opção)\\
    \opcao{Biblivre}\opcao{Koha}\opcao{Pergamum}\opcao{Sophia}\opcao{Aleph ou Alma}\opcao{Planilha ou registro manual}\opcao{Outro}
    \item Qual código de catalogação é adotado?\\
    \opcao{AACR2}\opcao{RDA}\opcao{Ambos (em transição)}\opcao{Não sei}
\end{enumerate}

\section*{Parte 3 --- Rotina de catalogação}

\begin{enumerate}
    \setcounter{enumi}{6}
    \item Quantos livros você cataloga, em média, por semana?\\
    \opcao{Nenhum}\opcao{1 a 10}\opcao{11 a 30}\opcao{31 a 100}\opcao{Mais de 100}
    \item Quanto tempo você leva, em média, para fazer a catalogação descritiva de um livro sem reaproveitar registros prontos?\\
    \opcao{Menos de 5 min}\opcao{5 a 10 min}\opcao{10 a 20 min}\opcao{20 a 30 min}\opcao{Mais de 30 min}
    \item Com que frequência você reaproveita registros de outras instituições (catalogação por cópia)?\\
    \opcao{Sempre}\opcao{Frequentemente}\opcao{Às vezes}\opcao{Raramente}\opcao{Nunca}
    \item O sistema da sua biblioteca permite importar registros MARC (arquivo \texttt{.mrc})?\\
    \opcao{Sim, e já utilizei}\opcao{Sim, mas nunca utilizei}\opcao{Não}\opcao{Não sei}
    \item Quais etapas da catalogação descritiva mais tomam o seu tempo? (marque até três)\\
    \opcao{Buscar os dados da obra}\opcao{Digitar os dados campo a campo}\opcao{Aplicar a pontuação}\opcao{Definir os pontos de acesso}\opcao{Preencher campos codificados (líder, 008)}\opcao{Revisar e corrigir erros}\opcao{Outra}
    \item Com que frequência você precisa corrigir erros de digitação ou de preenchimento em registros?\\
    \opcao{Nunca}\opcao{Raramente}\opcao{Às vezes}\opcao{Frequentemente}\opcao{Sempre}
\end{enumerate}

\section*{Parte 4 --- Funcionalidades da ferramenta}

O Cataloga+ é uma ferramenta web gratuita que, a partir do ISBN, busca os dados de um livro em fontes públicas, permite que o bibliotecário revise e corrija esses dados e gera o registro em MARC 21 (arquivo \texttt{.mrc}), pronto para ser importado no sistema da biblioteca. [Inserir link para vídeo de demonstração ou para a ferramenta.]

Considerando a sua prática de catalogação, indique o quanto cada funcionalidade ajudaria no seu dia a dia (1 = não ajudaria; 5 = ajudaria muito).

\begin{longtable}{|p{0.06\textwidth}|p{0.59\textwidth}|c|c|c|c|c|}
    \hline
    & \textbf{Funcionalidade} & \textbf{1} & \textbf{2} & \textbf{3} & \textbf{4} & \textbf{5} \\ \hline
    \endhead
    F1 & Busca automática dos dados do livro a partir do ISBN & \caixa & \caixa & \caixa & \caixa & \caixa \\ \hline
    F2 & Consulta a várias fontes ao mesmo tempo (CBL, Mercado Editorial, Open Library e Google Books), completando os dados que faltam em uma fonte com os de outra & \caixa & \caixa & \caixa & \caixa & \caixa \\ \hline
    F3 & Informar vários ISBNs de uma vez, digitando uma lista ou usando leitor de código de barras, e trabalhar em lote & \caixa & \caixa & \caixa & \caixa & \caixa \\ \hline
    F4 & Formulário para revisar e corrigir os dados antes de gerar o registro & \caixa & \caixa & \caixa & \caixa & \caixa \\ \hline
    F5 & Pré-visualização do registro MARC 21 antes de exportar & \caixa & \caixa & \caixa & \caixa & \caixa \\ \hline
    F6 & Geração do arquivo \texttt{.mrc} para importação no sistema da biblioteca, individualmente ou em lote & \caixa & \caixa & \caixa & \caixa & \caixa \\ \hline
    F7 & Escolha entre gerar o registro segundo o RDA ou o AACR2 & \caixa & \caixa & \caixa & \caixa & \caixa \\ \hline
    F8 & Aplicação automática da pontuação ISBD & \caixa & \caixa & \caixa & \caixa & \caixa \\ \hline
    F9 & Inversão automática dos nomes dos autores nos pontos de acesso (``Azevedo, Aluísio'') & \caixa & \caixa & \caixa & \caixa & \caixa \\ \hline
    F10 & Preenchimento automático de códigos, como país e idioma (008) e caracteres desconsiderados na ordenação de títulos com artigo (245) & \caixa & \caixa & \caixa & \caixa & \caixa \\ \hline
    F11 & Uso direto no navegador, sem instalação e sem cadastro & \caixa & \caixa & \caixa & \caixa & \caixa \\ \hline
    F12 & Dados do trabalho armazenados apenas no computador do usuário & \caixa & \caixa & \caixa & \caixa & \caixa \\ \hline
\end{longtable}

\section*{Parte 5 --- Percepção sobre a ferramenta}

Indique o seu grau de concordância com cada afirmação (1 = discordo totalmente; 5 = concordo totalmente).

\begin{longtable}{|p{0.06\textwidth}|p{0.59\textwidth}|c|c|c|c|c|}
    \hline
    & \textbf{Afirmação} & \textbf{1} & \textbf{2} & \textbf{3} & \textbf{4} & \textbf{5} \\ \hline
    \endhead
    \multicolumn{7}{|l|}{\textit{Utilidade percebida}} \\ \hline
    U1 & Uma ferramenta como esta reduziria o tempo que gasto catalogando livros & \caixa & \caixa & \caixa & \caixa & \caixa \\ \hline
    U2 & Uma ferramenta como esta reduziria erros de digitação e retrabalho & \caixa & \caixa & \caixa & \caixa & \caixa \\ \hline
    U3 & Uma ferramenta como esta contribuiria para a padronização dos registros & \caixa & \caixa & \caixa & \caixa & \caixa \\ \hline
    U4 & Com uma ferramenta como esta, eu poderia dedicar mais tempo a outras atividades, como o atendimento ao usuário & \caixa & \caixa & \caixa & \caixa & \caixa \\ \hline
    \multicolumn{7}{|l|}{\textit{Facilidade de uso percebida}} \\ \hline
    E1 & Acredito que aprenderia a usar esta ferramenta rapidamente & \caixa & \caixa & \caixa & \caixa & \caixa \\ \hline
    E2 & Seria fácil incorporar esta ferramenta à minha rotina sem trocar o sistema da biblioteca & \caixa & \caixa & \caixa & \caixa & \caixa \\ \hline
    \multicolumn{7}{|l|}{\textit{Confiança e intenção de uso}} \\ \hline
    C1 & Eu confiaria nos dados recuperados automaticamente, desde que pudesse revisá-los & \caixa & \caixa & \caixa & \caixa & \caixa \\ \hline
    I1 & Eu usaria uma ferramenta como esta na minha rotina de catalogação & \caixa & \caixa & \caixa & \caixa & \caixa \\ \hline
    I2 & Eu recomendaria esta ferramenta a colegas de profissão & \caixa & \caixa & \caixa & \caixa & \caixa \\ \hline
\end{longtable}

\section*{Parte 6 --- Questões abertas}

\begin{enumerate}
    \setcounter{enumi}{12}
    \item Qual das funcionalidades apresentadas seria mais útil para você? Por quê?
    \item O que poderia impedir você de utilizar uma ferramenta como esta? (pode marcar mais de uma opção)\\
    \opcao{Falta de confiança nos dados}\opcao{O sistema da biblioteca não importa MARC}\opcao{Política da instituição}\opcao{Acesso limitado à internet}\opcao{Prefiro catalogar manualmente}\opcao{Outro}
    \item Que outras fontes de dados ou funcionalidades você gostaria que a ferramenta tivesse?
    \item Deixe aqui comentários ou sugestões.
\end{enumerate}

% -----------------------------------------------------------------------------
\chapter{Questionário de avaliação de uso}
\label{ape:questionario-avaliacao}

Instrumento aplicado aos bibliotecários que participarem dos testes da ferramenta, imediatamente após a execução do roteiro de tarefas. Cada participante é identificado apenas por um código (P01, P02 etc.). Tempo estimado: 40 a 50 minutos, incluindo as tarefas.

\section*{Parte 1 --- Termo de consentimento e perfil}

Apresenta-se o mesmo termo de consentimento do \autoref{ape:questionario-adesao}, acrescido da informação de que o participante utilizará a ferramenta e de que o tempo de execução das tarefas será registrado pelo pesquisador. Em seguida, são repetidas as questões 1, 2, 5, 6 e 8 do \autoref{ape:questionario-adesao}, que caracterizam a formação, a experiência, o sistema utilizado, o código de catalogação e o tempo médio de catalogação manual do participante.

\section*{Parte 2 --- Roteiro de tarefas}

O participante recebe cinco livros físicos e uma lista com seus ISBNs, à qual é acrescentado um ISBN com dígito verificador incorreto. O pesquisador registra, para cada tarefa, se ela foi concluída sem ajuda (S), com ajuda (A) ou não concluída (N), e o tempo de execução. Ao fim de cada tarefa, o participante responde: ``De modo geral, esta tarefa foi:'', de 1 (muito difícil) a 7 (muito fácil).

\begin{longtable}{|p{0.05\textwidth}|p{0.48\textwidth}|p{0.11\textwidth}|p{0.09\textwidth}|p{0.12\textwidth}|}
    \hline
    \textbf{ID} & \textbf{Tarefa} & \textbf{Conclusão (S/A/N)} & \textbf{Tempo} & \textbf{Facilidade (1 a 7)} \\ \hline
    \endhead
    T1 & Buscar os dados de um livro pelo ISBN e abrir seu formulário de revisão & & & \\ \hline
    T2 & Buscar de uma só vez os demais ISBNs da lista & & & \\ \hline
    T3 & Revisar um registro, comparando-o com o livro em mãos, e corrigir os campos divergentes & & & \\ \hline
    T4 & Visualizar o registro MARC e localizar nele o título e os dados de publicação & & & \\ \hline
    T5 & Exportar o registro de um livro no padrão RDA & & & \\ \hline
    T6 & Identificar o documento com erro e tratá-lo (tentar novamente ou remover) & & & \\ \hline
    T7 & Alterar o padrão para AACR2 e exportar todo o lote & & & \\ \hline
    T8 & (Opcional) Importar o arquivo exportado no sistema da biblioteca & & & \\ \hline
\end{longtable}

Para cada livro, o pesquisador registra ainda o tempo desde a busca até a exportação e os campos alterados pelo participante durante a revisão:

\begin{longtable}{|p{0.18\textwidth}|p{0.58\textwidth}|p{0.12\textwidth}|}
    \hline
    \textbf{Livro (ISBN)} & \textbf{Campos corrigidos} & \textbf{Tempo} \\ \hline
    \endhead
    & \caixa~título \caixa~subtítulo \caixa~autores \caixa~editora \caixa~local \caixa~edição \caixa~ano \caixa~páginas \caixa~altura \caixa~formato \caixa~assuntos \caixa~sinopse \caixa~outro & \\ \hline
    & \textit{(repetir uma linha por livro)} & \\ \hline
\end{longtable}

\section*{Parte 3 --- Usabilidade (\textit{System Usability Scale})}

Indique o seu grau de concordância com cada afirmação (1 = discordo totalmente; 5 = concordo totalmente).

\begin{longtable}{|p{0.06\textwidth}|p{0.59\textwidth}|c|c|c|c|c|}
    \hline
    & \textbf{Afirmação} & \textbf{1} & \textbf{2} & \textbf{3} & \textbf{4} & \textbf{5} \\ \hline
    \endhead
    S1 & Eu acho que gostaria de usar esta ferramenta com frequência & \caixa & \caixa & \caixa & \caixa & \caixa \\ \hline
    S2 & Eu achei a ferramenta desnecessariamente complexa & \caixa & \caixa & \caixa & \caixa & \caixa \\ \hline
    S3 & Eu achei a ferramenta fácil de usar & \caixa & \caixa & \caixa & \caixa & \caixa \\ \hline
    S4 & Eu acho que precisaria do apoio de um técnico para conseguir usar esta ferramenta & \caixa & \caixa & \caixa & \caixa & \caixa \\ \hline
    S5 & Eu achei que as diversas funções da ferramenta estavam bem integradas & \caixa & \caixa & \caixa & \caixa & \caixa \\ \hline
    S6 & Eu achei que havia muita inconsistência na ferramenta & \caixa & \caixa & \caixa & \caixa & \caixa \\ \hline
    S7 & Eu imagino que a maioria das pessoas aprenderia a usar esta ferramenta rapidamente & \caixa & \caixa & \caixa & \caixa & \caixa \\ \hline
    S8 & Eu achei a ferramenta muito complicada de usar & \caixa & \caixa & \caixa & \caixa & \caixa \\ \hline
    S9 & Eu me senti muito confiante usando a ferramenta & \caixa & \caixa & \caixa & \caixa & \caixa \\ \hline
    S10 & Eu precisei aprender muitas coisas antes de conseguir usar esta ferramenta & \caixa & \caixa & \caixa & \caixa & \caixa \\ \hline
\end{longtable}

\section*{Parte 4 --- Funcionalidade e qualidade dos registros}

Indique o seu grau de concordância com cada afirmação (1 = discordo totalmente; 5 = concordo totalmente). Marque ``N/A'' quando a afirmação não se aplicar ao seu teste.

\begin{longtable}{|p{0.06\textwidth}|p{0.52\textwidth}|c|c|c|c|c|c|}
    \hline
    & \textbf{Afirmação} & \textbf{1} & \textbf{2} & \textbf{3} & \textbf{4} & \textbf{5} & \textbf{N/A} \\ \hline
    \endhead
    Q1 & Os dados recuperados estavam corretos na maior parte dos campos & \caixa & \caixa & \caixa & \caixa & \caixa & \caixa \\ \hline
    Q2 & Os dados recuperados foram suficientes para a catalogação descritiva das obras & \caixa & \caixa & \caixa & \caixa & \caixa & \caixa \\ \hline
    Q3 & A combinação das fontes preencheu campos que eu teria de digitar & \caixa & \caixa & \caixa & \caixa & \caixa & \caixa \\ \hline
    Q4 & O registro gerado está de acordo com o código escolhido (RDA ou AACR2) & \caixa & \caixa & \caixa & \caixa & \caixa & \caixa \\ \hline
    Q5 & A pontuação ISBD foi aplicada corretamente & \caixa & \caixa & \caixa & \caixa & \caixa & \caixa \\ \hline
    Q6 & Os pontos de acesso de autor (nomes invertidos) estavam corretos & \caixa & \caixa & \caixa & \caixa & \caixa & \caixa \\ \hline
    Q7 & A pré-visualização facilitou a conferência do registro & \caixa & \caixa & \caixa & \caixa & \caixa & \caixa \\ \hline
    Q8 & O arquivo exportado foi importado no sistema da biblioteca sem erros & \caixa & \caixa & \caixa & \caixa & \caixa & \caixa \\ \hline
    Q9 & Os dados exibidos no sistema da biblioteca corresponderam aos que revisei & \caixa & \caixa & \caixa & \caixa & \caixa & \caixa \\ \hline
    Q10 & As mensagens de erro deixaram claro o que havia acontecido & \caixa & \caixa & \caixa & \caixa & \caixa & \caixa \\ \hline
    Q11 & O trabalho em lote é adequado ao volume de obras que catalogo & \caixa & \caixa & \caixa & \caixa & \caixa & \caixa \\ \hline
\end{longtable}

\section*{Parte 5 --- Eficiência percebida}

\begin{enumerate}
    \item Em comparação com a catalogação manual, catalogar com o Cataloga+ foi:\\
    \opcao{Muito mais lento}\opcao{Mais lento}\opcao{Igual}\opcao{Mais rápido}\opcao{Muito mais rápido}
    \item Quanto tempo, em minutos, você estima levar para catalogar um livro manualmente? \underline{\hspace{1.5cm}} E com o Cataloga+? \underline{\hspace{1.5cm}}
    \item A ferramenta reduziu a quantidade de dados que precisei digitar.\\
    (1 = discordo totalmente; 5 = concordo totalmente) \quad 1~\caixa\quad 2~\caixa\quad 3~\caixa\quad 4~\caixa\quad 5~\caixa
\end{enumerate}

\section*{Parte 6 --- Satisfação e intenção de uso}

\begin{enumerate}
    \setcounter{enumi}{3}
    \item Eu usaria o Cataloga+ na minha rotina de catalogação.\\
    (1 = discordo totalmente; 5 = concordo totalmente) \quad 1~\caixa\quad 2~\caixa\quad 3~\caixa\quad 4~\caixa\quad 5~\caixa
    \item Em uma escala de 0 a 10, quanto você recomendaria o Cataloga+ a um colega de profissão? \underline{\hspace{1.5cm}}
\end{enumerate}

\section*{Parte 7 --- Questões abertas}

\begin{enumerate}
    \setcounter{enumi}{5}
    \item Quais foram os pontos positivos da ferramenta?
    \item Quais dificuldades você encontrou?
    \item Você identificou algum erro nos dados ou nos registros gerados? Descreva.
    \item Que campos, regras de catalogação ou funcionalidades fizeram falta?
    \item Deixe aqui outros comentários ou sugestões.
\end{enumerate}
